% =====================================================================
% REPLACEMENT for supplementary_results.tex section
%   "Climate is not a robust driver of sorting or diagnosticity" (\label{sr:climate})
% Based on the completed genome-wide, eMLG-level Module C analysis.
% =====================================================================

\section{Climate association does not explain repeated sorting or diagnostic enrichment}\label{sr:climate}

We asked whether climate-association evidence was distributed non-randomly with respect
to parental diagnosticity, recombination rate, or repeated directional sorting. The
analysis used the complete set of 32{,}840 eMLGs from the primary {\AA}land-excluded
BayPass configuration. For each climate axis and each of 10{,}000 null covariates drawn
from the fixed among-population covariance matrix, the genome-wide Bayes-factor vector
was reduced to the same three rank-based statistics. This calibration retains the
population structure represented by $\Omega$ and avoids defining climate-associated
regions by an arbitrary Bayes-factor threshold.

\paragraph{Directional sorting and recombination.} Neither directional sorting
nor recombination rate was associated with climate-association strength beyond the
structured null. Among parentally differentiated eMLGs, the mean Bayes-factor-percentile
difference between directional and non-directional units was 0.055 for PC1 (empirical
$P=0.107$, FDR $=0.320$) and 0.024 for PC2 ($P=0.475$, FDR $=0.571$), compared with a
null median of 0.033 and a central 95\% interval of 0.007--0.057. Correlations with
recombination were likewise null (PC1 $\rho=0.001$, FDR $=0.909$; PC2 $\rho=-0.006$,
FDR $=0.571$), as were correlations with continuous sorting magnitude (PC1 $P=0.575$;
PC2 $P=0.357$).

\paragraph{Diagnostic content.} Diagnostic Index was unrelated to climate-association
strength on PC1 ($\rho=0.017$, FDR $=0.571$). PC2 showed the only departure among the
six primary tests: its Bayes-factor rank was negatively correlated with DI
($\rho=-0.125$, empirical $P=0.001$, FDR $=0.006$), below the null 95\% interval of
$-0.093$ to 0.047. Because higher DI denotes greater parental diagnosticity, this weak
correlation does not indicate enrichment among the most diagnostic eMLGs. The raw-Bayes-
factor correlation was similar ($r=-0.119$, $P<0.001$), but the effect was diffuse: no
individual eMLG survived the locus-level simulated-FDR analysis. It therefore describes
a genome-wide shift in ranks, not a set of supported climate-adaptation loci.

\paragraph{Sensitivity analyses and validation.} Upper-tail summaries were descriptive
because extreme-rank membership was more sensitive to BayPass Monte Carlo variation. The
top PC1 fractions contained more directionally sorted eMLGs than expected under the null
(53.8\%, 47.8\%, and 39.6\% in the upper 0.1\%, 0.5\%, and 1\%; empirical $P=0.028$,
0.0008, and 0.0034) and tended toward lower recombination. These cut-off-dependent
patterns were absent for PC2 and did not alter the primary tests. Regenerated and archived
per-eMLG exceedance counts correlated at 0.99996 and 0.99995 for PC1 and PC2, with total
count differences of 0.0019\% and 0.0046\%. Independent runs showed that Monte Carlo
variation in each primary statistic was at most 6.2\% of its variation among null
covariates. Overall, the measured climate gradients do not robustly organise the repeated
direction or magnitude of sorting. The weak PC2--DI relationship remains conditional and
exploratory because climate, ancestry, and population history are geographically
entangled.
